\begin{ouverture}[college3]{11}{Énergie cinétique et bilan mécanique}{Partie 3 · L’énergie, ses transferts et ses conversions}{Pourquoi la vitesse pèse-t-elle autant sur un freinage ?}
\objectif{Calculer une énergie cinétique en unités SI.}
\objectif{Comparer l’effet de la masse et de la vitesse.}
\objectif{Suivre une conversion et distinguer énergie et force.}
\prerequis{Vitesse en m/s ; carré d’un nombre.}
\prerequis{Conservation et formes d’énergie.}
\end{ouverture}
\pcPage{Le cours}{Énergie cinétique et bilan mécanique}
\phantomsection\label{ess-11}
\essSection{1}{L’énergie du mouvement}
Un objet de masse m et de vitesse de valeur v possède, dans le référentiel choisi, une énergie cinétique \[E_c=\dfrac12 mv^2.\] Avec m en kg et v en m/s, $E_c$ est en joules, J. La vitesse se convertit avant le carré : une valeur en km/h doit être divisée par 3,6. Un objet immobile dans le référentiel considéré y a une énergie cinétique nulle.\par
\pcFigure{figures/b11-ec.pdf}{À masse fixée, le graphe d’Ec en fonction de v n’est pas une droite.}
\essSection{2}{Des facteurs à comparer}
À vitesse fixée, doubler la masse double l’énergie cinétique. À masse fixée, doubler la vitesse multiplie cette énergie par quatre ; tripler la vitesse la multiplie par neuf. Ces effets expliquent pourquoi une petite augmentation de vitesse peut avoir un effet énergétique important. Le calcul ne donne pas à lui seul une distance d’arrêt : freinage, adhérence, pente et temps de réaction interviennent aussi.\par
\clearpage
\essSection{3}{Position et conversion}
Une énergie potentielle de pesanteur dépend de la position et du choix d’un niveau de référence. Lors d’une chute, elle peut diminuer tandis que l’énergie cinétique augmente. Si les frottements sont négligés, leur somme, l’énergie mécanique, reste constante. Avec frottements, une partie de l’énergie mécanique est transférée ou convertie en énergie thermique ; l’énergie totale d’un système suffisamment large reste conservée.\par
\pcFigure{figures/b11-chute.pdf}{Modèle de chute : exemple chiffré sans frottement ; le niveau de référence est annoncé.}
\essSection{4}{Freiner et argumenter}
Freiner réduit l’énergie cinétique et augmente notamment l’énergie thermique des freins, pneus et milieu. L’énergie n’est pas détruite. Une force se mesure en N, une énergie en J, une puissance en W : ce sont des grandeurs distinctes. Au Brevet, comparer les carrés des vitesses peut suffire à justifier un rapport d’énergies sans calculer deux grandes valeurs.\par
\begin{pcEncadre}{method}{Un exemple expliqué}Pour un ensemble vélo-cycliste de 80 kg à 18 km/h, $v=18/3{,}6=5\unite{m/s}$ ; $E_c=\tfrac12\times80\times5^2=1\,000\unite J$. À 36 km/h, v double : $E_c=4\,000\unite J$. La comparaison suppose la même masse et le même référentiel.\end{pcEncadre}
\pcPage{Expérimenter et retenir}{Énergie cinétique et bilan mécanique}
\phantomsection\label{activite-11}
\begin{pcEncadre}{method}{Investigation en classe}Au tableur, comparer les énergies d’une même masse pour 2, 4, 6 et 8 m/s. Ajouter une colonne $v^2$ et tracer Ec en fonction de v, puis de $v^2$. Exploiter ensuite une vidéo ou une simulation de chute en distinguant modèle sans frottement et observation réelle. Ne pas provoquer de choc dangereux.\end{pcEncadre}
\pcSousTitre{Ce que je dois savoir expliquer}\begin{questions}
\item Écrire Ec avec m en kg et v en m/s.
\item Justifier l’effet du carré de la vitesse.
\item Faire un bilan comprenant les sorties thermiques.
\end{questions}
\begin{pcEncadre}{attention}{Une erreur à éviter}Doubler la vitesse ne double pas l’énergie cinétique : dans le modèle, cela la multiplie par quatre. Il faut mettre toute la vitesse convertie au carré.\end{pcEncadre}
\pcRepere{cor-activite-11}{Pistes pour l’investigation}\pcGrille
\clearpage\markright{Je démarre}
\rubrique{Je démarre}
\par\noindent\begin{minipage}{\linewidth}
\begin{exercice}[1]{La forme d’énergie}
\begin{questions}
\item Quelle énergie est associée au mouvement d’un objet ?
\item Dans quelle unité SI s’exprime-t-elle ?
\end{questions}
\end{exercice}
\end{minipage}
\par\noindent\begin{minipage}{\linewidth}
\begin{exercice}[1]{Les bonnes unités}
\begin{questions}
\item Quelles unités de m et v employer dans Ec = ½mv² ?
\item Convertir 500 g en kg.
\end{questions}
\end{exercice}
\end{minipage}
\par\noindent\begin{minipage}{\linewidth}
\begin{exercice}[1]{La vitesse}
\begin{questions}
\item Convertir 36 km/h en m/s.
\item Convertir 72 km/h en m/s.
\end{questions}
\end{exercice}
\end{minipage}
\par\noindent\begin{minipage}{\linewidth}
\begin{exercice}[1]{Un calcul direct}
\begin{questions}
\item Calculer Ec pour 2 kg à 3 m/s.
\item Que vaut Ec si cet objet est immobile dans le référentiel ?
\end{questions}
\end{exercice}
\end{minipage}
\par\noindent\begin{minipage}{\linewidth}
\begin{exercice}[1]{Deux facteurs}
\begin{questions}
\item À vitesse identique, la masse double. Que devient Ec ?
\item À masse identique, la vitesse double ?
\end{questions}
\end{exercice}
\end{minipage}
\clearpage\markright{Je m’entraîne}
\rubrique{Je m’entraîne}
\par\noindent\begin{minipage}{\linewidth}
\begin{exercice}[2]{Le ballon}
\begin{questions}
\item Un ballon de 400 g va à 5 m/s. Calculer Ec.
\item Si v devient 10 m/s, calculer la nouvelle valeur.
\end{questions}
\end{exercice}
\end{minipage}
\par\noindent\begin{minipage}{\linewidth}
\begin{exercice}[2]{Une voiture}
\begin{questions}
\item Une voiture de 1 000 kg va à 54 km/h. Déterminer v puis Ec.
\item Donner Ec en kJ.
\end{questions}
\end{exercice}
\end{minipage}
\par\noindent\begin{minipage}{\linewidth}
\begin{exercice}[2]{Ralentir}
\begin{questions}
\item Un ensemble de 60 kg passe de 10 à 5 m/s. Calculer les deux énergies.
\item Quelle énergie cinétique est perdue ?
\end{questions}
\end{exercice}
\end{minipage}
\par\noindent\begin{minipage}{\linewidth}
\begin{exercice}[2]{Une chute idéale}
\begin{questions}
\item Au départ, Ec = 0 J et Ep = 40 J. Sans frottement, que vaut Ec lorsque Ep = 15 J ?
\item Pourquoi faut-il annoncer le modèle ?
\end{questions}
\end{exercice}
\end{minipage}
\par\noindent\begin{minipage}{\linewidth}
\begin{exercice}[2]{Le rapport}
\begin{questions}
\item À masse fixée, comparer Ec à 30 m/s et à 10 m/s.
\item À vitesse fixée, comparer 120 kg et 80 kg.
\end{questions}
\end{exercice}
\end{minipage}
\clearpage\markright{Je raisonne}
\rubrique{Je raisonne}
\par\noindent\begin{minipage}{\linewidth}
\begin{exercice}[3]{Une hausse de 20 \%}
\begin{questions}
\item La vitesse est multipliée par 1,2 à masse fixée. Par quel facteur Ec est-elle multipliée ?
\item Exprimer l’augmentation en pourcentage.
\end{questions}
\end{exercice}
\end{minipage}
\par\noindent\begin{minipage}{\linewidth}
\begin{exercice}[3]{Changer de référentiel}
\begin{questions}
\item Un passager est immobile par rapport à son train. Peut-il avoir une énergie cinétique dans le référentiel du sol ?
\item Quelle précision manque à « cet objet n’a aucune énergie cinétique » ?
\end{questions}
\end{exercice}
\end{minipage}
\par\noindent\begin{minipage}{\linewidth}
\begin{exercice}[3]{La force n’est pas l’énergie}
\begin{questions}
\item Un élève donne une énergie cinétique de 250 N. Corriger l’unité.
\item Peut-on trouver une distance de freinage avec seulement Ec ?
\end{questions}
\end{exercice}
\end{minipage}
\par\noindent\begin{minipage}{\linewidth}
\begin{exercice}[3]{Un bilan réel}
\begin{questions}
\item Une chute commence avec 100 J d’énergie mécanique et finit avec 85 J. Où les 15 J peuvent-ils être passés ?
\item La conservation totale est-elle remise en cause ?
\end{questions}
\end{exercice}
\end{minipage}
\par\noindent\begin{minipage}{\linewidth}
\begin{exercice}[3]{Comparer deux véhicules}
\begin{questions}
\item A : 800 kg à 20 m/s ; B : 1 600 kg à 10 m/s. Lequel a le plus d’énergie cinétique ?
\item Pourquoi comparer seulement les masses donnerait-il une erreur ?
\end{questions}
\end{exercice}
\end{minipage}
\pcPage{Défi du chapitre}{Énergie cinétique et bilan mécanique}
\phantomsection\label{defi-ex-11}
\pcDefiTitre{11}{Une étude de sécurité}
\begin{questions}
\item Une voiture et ses occupants ont une masse de 1 200 kg. Comparer Ec à 50 km/h et 100 km/h.
\item Expliquer le résultat sans recalculer la seconde énergie.
\item Rédiger une conclusion de quatre lignes qui distingue énergie, distance d’arrêt et conditions de freinage.
\end{questions}
\pcGrille
